\section*{第 \ref{ch:b2:metamorphic-facies} 章——变质反应与变质相}

\begin{solution}{exo:b2:metamorphic-facies:1}
(a) 接触\omterm{def:b2:metamorphic-facies:metamorphism}{变质作用}，即由侵入体热量引起的局地热\omterm{def:b2:metamorphic-facies:metamorphism}{变质作用}。
(b) 位错\omterm{def:b2:metamorphic-facies:metamorphism}{变质作用}，发生在断层带中：粒度减小，形成糜棱岩。
(c) 洋底\omterm{def:b2:metamorphic-facies:metamorphism}{变质作用}，由洋中脊附近的热海水引起，
重结晶不完全。(d) 埋藏\omterm{def:b2:metamorphic-facies:metamorphism}{变质作用}，区域性，未变形，
亦无片理。
\end{solution}

\begin{solution}{exo:b2:metamorphic-facies:2}
$\Delta_r S = 95.40 - 82.80 = \qty{12.60}{J.K^{-1}}$，$\Delta_r V = 4.986 - 4.415 =
\qty{0.571}{J.bar^{-1}}$；斜率 $12.60/0.571 = \qty{22.1}{bar/K} = \qty{2.21}{MPa/K}$。
温度升高 \qty{100}{\celsius} 时，平衡压强升高 \qty{2.21}{kbar}。
\end{solution}

\begin{solution}{exo:b2:metamorphic-facies:3}
(a) 绿片岩相。(b) 蓝片岩相。(c) 角闪岩相：角闪石与斜长石，
其钙长石含量超过 \qty{17}{\percent}。(d) 榴辉岩相。
\end{solution}

\begin{solution}{exo:b2:metamorphic-facies:4}
(a) $\rho = M/V$：蓝晶石 $162.046/44.15 = \qty{3.670}{g.cm^{-3}}$，红柱石
$162.046/51.52 = 3.145$，夕线石 $162.046/49.86 = 3.250$。(b) 蓝晶石：
压强较高时，吉布斯自由能中的项 $V\,(P - 1)$ 对体积最小者增长最慢。
(c) 天然晶体含有其他元素（例如三价铁替代了天然红柱石和夕线石中的部分
铝），且天然样品的测量
精度低于表中给出的摩尔体积。
\end{solution}

\begin{solution}{exo:b2:metamorphic-facies:5}
(a) $\Delta_r G = 3900 - 298.15\times8.59 = \qty{+1339}{J} > 0$：蓝晶石。(b) 在
\qty{1}{bar} 下，$\Delta_r H = T\Delta_r S$：$T = 3900/8.59 = \qty{454}{K} =
\qty{181}{\celsius}$。(c) $P = (8.59\times673.15 - 3900)/0.737 = 1882/0.737 =
\qty{2554}{bar} = \qty{2.55}{kbar}$，与实测值 2.54 相比。(d) 低温下反应
太慢而无法进行：红柱石以亚稳态保留，与大多数高温
矿物在地表的情形相同。
\end{solution}

\begin{solution}{exo:b2:metamorphic-facies:6}
(a) $\Delta_r H = 7700 + 1000\times(121.39 - 122.29) = \qty{6800}{J}$；$\Delta_r S = 296.44
- 285.02 = \qty{11.42}{J.K^{-1}}$；$\Delta_r G = 6800 - 1000\times11.42 = \qty{-4620}{J}$。
(b) 取 \qty{298.15}{K} 下的数值：$7700 - 1000\times12.60 = \qty{-4900}{J}$。
(c) 夕线石。简化方法的误差为 \qty{280}{J}，约占 \qty{6}{\percent}；除以
$\Delta_r V = 0.571$，沿该线相当于约 \qty{490}{bar} 的压强，远小于
焓的不确定度所带来的影响。
\end{solution}

\begin{solution}{exo:b2:metamorphic-facies:7}
(a) $\Delta_r G = -1127.4 + 1128.5 = \qty{1.1}{kJ}$，且 $\Delta_r V = (34.15 -
36.93)/10 = \qty{-0.278}{J.bar^{-1}}$：$P = 1100/0.278 = \qty{3957}{bar} \approx
\qty{4.0}{kbar}$。(b) $\Delta_r S = 88.0 - 91.7 = -3.7$；斜率 $-3.7/-0.278 =
\qty{13.3}{bar/K}$。(c) 令 $\Delta_r H = 0$，$P = T\Delta_r S/\Delta_r V = 13.31\,T$：
在 \qty{573.15}{K} 时，\qty{7629}{bar} $= \qty{7.63}{kbar}$。(d) 贝壳中的文石
在地表呈亚稳态：生物体形成了它，它之所以能保存下来，是因为它
只能缓慢地转变为方解石。
\end{solution}

\begin{solution}{exo:b2:metamorphic-facies:8}
(a) $-32.4/-1.698 = \qty{19.1}{bar/K}$。(b) $P = (T\Delta_r S - \Delta_r H)/\Delta_r V =
(-32.4\times773.15 + 5000)/(-1.698) = (-25050 + 5000)/(-1.698) = \qty{11808}{bar}
\approx \qty{11.8}{kbar}$。(c) $\sqrt{2.6^2 + 3.6^2 + 1.0^2} = \qty{4.55}{kJ}$，即作用于
$\Delta_r H$，故压强变化 $4550/1.698 = \qty{2.7}{kbar}$。(d) 在
榴辉岩相的压强下，斜长石中的钠长石组分相对于
硬玉和石英是不稳定的；钠进入辉石（绿辉石）中，斜长石
消失。
\end{solution}

\begin{solution}{exo:b2:metamorphic-facies:9}
(a) 从 1 到 \qty{2}{kbar}：$1000/34 = \qty{29.4}{bar/K}$；从 2 到 3：$1000/25 =
\qty{40.0}{bar/K}$。(b) 在 $\Delta_r S$ 不变时，$\Delta_r V$ 与
斜率成反比：它缩小的倍数为 $40.0/29.4 = 1.36$。(c) 释放出的水：其摩尔
体积随压强升高而减小，而固体的体积几乎不变。
\end{solution}

\begin{solution}{exo:b2:metamorphic-facies:10}
(a) 在 \qty{873.15}{K} 下：蓝晶石 $=$ 红柱石 $(8.59\times873.15 - 3900)/0.737 =
\qty{4885}{bar}$；蓝晶石 $=$ 夕线石 $(12.60\times873.15 - 7700)/0.571 =
\qty{5782}{bar}$；红柱石 $=$ 夕线石 $(4.01\times873.15 - 3800)/(-0.166) =
\qty{1799}{bar}$。(b) 在 \qty{4885}{bar} 处，低于蓝晶石 $=$ 夕线石线（5782），
夕线石比蓝晶石更稳定；压强高于红柱石 $=$ 夕线石线
（1799）时，夕线石比红柱石更稳定：该点位于夕线石域内，
蓝晶石 $=$ 红柱石平衡在此处为亚稳。(c) 在 \qty{4}{kbar} 时：介于 1799
与 \qty{5782}{bar} 之间为夕线石。在 \qty{7}{kbar} 时：高于 \qty{5782}{bar} 时蓝晶石优于
夕线石，高于 \qty{4885}{bar} 时又优于红柱石：蓝晶石。
\end{solution}

\begin{solution}{exo:b2:metamorphic-facies:11}
(a) $z = P/(\rho g) = 3.75\times10^8/(2700\times9.81) = \qty{14.2}{km}$，不确定度为
$0.25\times10^8/26487 = \qty{0.94}{km}$。(b) $(504 - 10)/14.16 =
\qty{35}{\celsius.km^{-1}}$。(c) 三种矿物仅在\omterm{def:b2:metamorphic-facies:invariant}{不变点}处才能平衡共存：
岩石的演化路径必须经过该点，且反应必须在该处
终止。大多数岩石从旁边经过，并保留一种或两种同质多象变体。
\end{solution}

\begin{solution}{exo:b2:metamorphic-facies:12}
(a) X：$T = 10 + 10\times30 = \qty{310}{\celsius}$（\qty{583.15}{K}）；$P = 3000\times
9.81\times30000 = \qty{8.83e8}{Pa} = \qty{8.83}{kbar}$。Y：$T = 10 + 30\times20 =
\qty{610}{\celsius}$（\qty{883.15}{K}）；$P = 2800\times9.81\times20000 = \qty{5.49}{kbar}$。
(b) X：方解石 $=$ 文石，在 $13.31\times583.15 = \qty{7.76}{kbar} < 8.83$ 时为文石；
钠长石 $=$ 硬玉 + 石英，在 $(32.4\times583.15 - 5000)/1.698 = \qty{8.18}{kbar} <
8.83$，硬玉 + 石英；蓝晶石 $=$ 红柱石，在 $(8.59\times583.15 - 3900)/0.737 =
\qty{1.51}{kbar}$，低于三相点温度：蓝晶石。Y：方解石 $=$
文石，在 \qty{11.75}{kbar} $> 5.49$ 时为方解石；钠长石 $=$ 硬玉 + 石英，在
\qty{13.9}{kbar} 时为钠长石；高于三相点时，蓝晶石 $=$ 夕线石，在
\qty{6.00}{kbar} $> 5.49$ 时，红柱石 $=$ 夕线石在 \qty{1.56}{kbar} $< 5.49$ 时：夕线石。
(c) X 对应 \qty{10}{\celsius\per\km} 的地温梯度，是较冷的路径：硬玉与
文石在仅 \qty{310}{\celsius} 时即将其归入蓝片岩相。Y 的地温梯度是 X 的三倍，
最终在 \qty{610}{\celsius} 与中等压强下形成夕线石，
属于角闪岩相。
\end{solution}

\begin{solution}{pb:b2:metamorphic-facies:1}
\noindent\textbf{第一部分。}

\textbf{1.} $\Delta_r H = -2589.9 + 2593.8 = \qty{3.9}{kJ}$; $\Delta_r S = 91.39 - 82.80 =
\qty{8.59}{J.K^{-1}}$; $\Delta_r V = 5.152 - 4.415 = \qty{0.737}{J.bar^{-1}}$.

\textbf{2.} $\qty{7.7}{kJ}$; $\qty{12.60}{J.K^{-1}}$; $\qty{0.571}{J.bar^{-1}}$.

\textbf{3.} $\qty{3.8}{kJ}$；$\qty{4.01}{J.K^{-1}}$；$\qty{-0.166}{J.bar^{-1}}$；每项
为第二者减去第一者，因红柱石相互抵消。

\textbf{4.} 夕线石的熵最大（95.40）：它位于其两条线的
高温一侧。蓝晶石的体积最小（4.415）：它位于其两条线的
高压一侧。体积最大的红柱石占据
低压一侧。

\textbf{5.} $8.59/0.737 = \qty{11.66}{bar/K}$; $12.60/0.571 = \qty{22.07}{bar/K}$;
$4.01/(-0.166) = \qty{-24.16}{bar/K}$.

\textbf{6.} 夕线石的熵比红柱石大而体积更小：$\Delta_r S > 0$
且 $\Delta_r V < 0$。升高温度有利于夕线石，升高压强
同样如此，故该线随温度升高指向更低压强。

\noindent\textbf{第二部分。}

\textbf{7.} $3900 - 298.15\times8.59 = \qty{+1339}{J}$; $7700 - 298.15\times12.60 =
\qty{+3943}{J}$。两者均为正，蓝晶石在地表稳定。

\textbf{8.} $T = \Delta_r H/\Delta_r S$: $3900/8.59 = \qty{454.0}{K} = \qty{181}{\celsius}$;
$7700/12.60 = \qty{611.1}{K} = \qty{338}{\celsius}$; $3800/4.01 = \qty{947.6}{K} =
\qty{674}{\celsius}$.

\textbf{9.} 在 \qty{1}{bar} 下加热时，蓝晶石在
\qty{181}{\celsius} 时转变为红柱石；在 \qty{338}{\celsius} 时蓝晶石已不再稳定，故蓝晶石 $=$
夕线石在此处为亚稳态；红柱石在
\qty{674}{\celsius} 时转变为夕线石。第一与第三个为稳定平衡。

\textbf{10.} $(8.59\times673.15 - 3900)/0.737 = \qty{2554}{bar} = \qty{2.55}{kbar}$,
与实验值相差 \qty{0.01}{kbar}。

\textbf{11.} $(12.60\times973.15 - 7700)/0.571 = 4562/0.571 = \qty{7989}{bar} =
\qty{7.99}{kbar}$，高于实验值 \qty{0.13}{kbar}。

\textbf{12.} 蓝晶石 $=$ 红柱石的线在 \qty{1}{bar} 下自 \qty{181}{\celsius} 起上升，斜率为
\qty{11.7}{bar/K}；红柱石 $=$ 夕线石的线自约 \qty{11.5}{kbar}（\qty{200}{\celsius}）
降至 \qty{1}{bar}（\qty{674}{\celsius}）；蓝晶石 $=$ 夕线石，两条上升线中较陡者，
在约 \qty{514}{\celsius} 与 \qty{3.9}{kbar} 处与两线相交，每条线仅在交点的
一侧稳定（本章图）。

\noindent\textbf{第三部分。}

\textbf{13.} $a_1 = 8.59/0.737 = \qty{11.655}{bar/K}$, $b_1 = 3900/0.737 =
\qty{5291.7}{bar}$; $a_2 = 12.60/0.571 = \qty{22.067}{bar/K}$, $b_2 = 7700/0.571 =
\qty{13485.1}{bar}$.

\textbf{14.} $T = (b_2 - b_1)/(a_2 - a_1) = 8193.4/10.412 = \qty{786.92}{K} =
\qty{513.8}{\celsius}$.

\textbf{15.} $P = 11.655\times786.92 - 5291.7 = \qty{3880}{bar} = \qty{3.88}{kbar}$.

\textbf{16.} $(4.01\times786.92 - 3800)/(-0.166) = (3155.5 - 3800)/(-0.166) =
\qty{3882}{bar}$：在数据舍入误差范围内二者压强相同，符合本章命题的要求。

\textbf{17.} \qty{514}{\celsius} 与 \qty{504}{\celsius} 相差 \qty{10}{\celsius}，在
\qty{20}{\celsius} 的不确定度之内；\qty{3.88}{kbar} 与 3.75 相差 \qty{0.13}{kbar}，
在 0.25 之内。吻合良好。

\textbf{18.} 否：\qty{536}{\celsius} 与实验值相差 \qty{32}{\celsius}，超出其
不确定度。忽略热容并非主要误差来源，主要误差来自焓。

\noindent\textbf{第四部分。}

\textbf{19.} $\Delta_r H$（蓝晶石 $\to$ 红柱石）增加 \qty{1000}{J}；在固定
$T$ 下，$P = (T\Delta_r S - \Delta_r H)/\Delta_r V$ 下降 $1000/0.737 = \qty{1357}{bar}$。

\textbf{20.} $\Delta T = -1357/(a_2 - a_1) = -1357/10.41 = \qty{-130}{K}$：三相点
移至约 \qty{384}{\celsius}。

\textbf{21.} 蓝晶石 $=$ 夕线石的线下降 $1000/0.571 = \qty{1751}{bar}$，而
三相点移动 $+1751/10.41 = \qty{+168}{K}$。

\textbf{22.} 单独将任一焓值改变 $\pm\qty{2.0}{kJ}$，三相点便
移动 $\pm\qty{260}{K}$（红柱石）或 $\pm\qty{336}{K}$（夕线石）：仅量热数据
即可允许数百度范围内的三相点。

\textbf{23.} 表中所列焓值并非纯粹由量热法得到：在其
不确定度范围内经过调整，使计算所得的线与逆转实验相符。
实验确定了相图；数据表则将其一致地推广到其他反应。

\textbf{24.} 由表中数据在性质取为常数时算得的三相点位于
\textbf{\qty[text-series-to-math=true]{514}{\celsius}}.
\end{solution}
